\documentclass[main]{subfiles}

\begin{document}
% Integration note: this file is not yet wired into main.tex. To include it,
% add \appendix immediately before its \subfile{06_appendix_notation} line
% (after \subfile{05_gap_analysis}), which switches chapter numbering to
% "Appendix A" for this and any subsequent appendix files. The
% \ifSubfilesClassLoaded guard below only fires when this file is compiled
% on its own, so it does not interfere with that.
\ifSubfilesClassLoaded{\appendix\setcounter{chapter}{0}}{}
\chapter{Notation and Symbol Reference}
\label{app:notation}

This appendix collects, in one place, every symbol, operator, and sign
convention used across Chapters~1--5. Where a quantity is introduced at
different levels of formality in different chapters (for example, the
free-particle versus full Foldy--Wouthuysen Hamiltonian), both forms are
given together with a note on how they relate.

\section{Dirac Algebra}

Table~\ref{tab:app-dirac} lists the gamma-matrix and Pauli-matrix notation
used throughout the Foldy--Wouthuysen derivations of Chapters~2--3, all in
the Dirac/Pauli representation with metric signature $(+,-,-,-)$.

\begin{table}[h]
\centering
\begin{tabular}{@{}ll>{\raggedright\arraybackslash}p{6cm}@{}}
\toprule
Symbol & Meaning & Context \\
\midrule
$\gamma^\mu$ & Dirac gamma matrices ($\mu=0,1,2,3$) & Relativistic fermion field theory \\
$\gamma^0$ & Time-like gamma matrix & Block-diagonal: $\mathrm{diag}(1,1,-1,-1)$ in the Dirac representation \\
$\gamma_5$ & $i\gamma^0\gamma^1\gamma^2\gamma^3$ & Chiral matrix; off-diagonal in the Dirac representation \\
$\beta$ & $\gamma^0$ & Used throughout the Foldy--Wouthuysen expansion \\
$\alpha^i$ & $\gamma^0\gamma^i$ ($i=1,2,3$) & Velocity operator; off-diagonal in the Dirac representation \\
$\Sigma^i$ & $\mathrm{diag}(\sigma^i,\sigma^i)$ & $4\times4$ spin matrix \\
$\sigma^i$ & Pauli matrices ($i=1,2,3$) & $2\times2$ spin operators, non-relativistic limit \\
$\boldsymbol\sigma$ & $(\sigma^1,\sigma^2,\sigma^3)$ & Pauli vector \\
$D_\mu$ & Gauge-covariant derivative & Appears in the general SME Lagrangian (Ch.~1, Ch.~2) \\
$\eta_\Gamma$ & CPT parity of a bilinear $\bar\psi\Gamma\psi$ & $+1$ for $\Gamma\in\{\mathbb 1,\gamma^\mu,\sigma^{\mu\nu}\}$, $-1$ for $\Gamma\in\{\gamma_5,\gamma_5\gamma^\mu\}$ (Ch.~2). Does not apply directly to derivative-coupled coefficients; see Section~\ref{app:cpt-table} \\
\bottomrule
\end{tabular}
\caption{Dirac algebra notation used in Chapters~2--3.}
\label{tab:app-dirac}
\end{table}

\section{SME Coefficients}

Table~\ref{tab:app-sme} lists every SME fermion-sector coefficient discussed
in the thesis, with its CPT and Lorentz parity and its leading-order
non-relativistic role.

\begin{table}[h]
\centering
\small
\begin{tabular}{@{}llcc>{\raggedright\arraybackslash}p{5.5cm}@{}}
\toprule
Symbol & Type & CPT & Lorentz & Definition / leading role \\
\midrule
$a_\mu$ & Vector & Odd & Odd & $\mathcal L_a=-a_\mu\bar\psi\gamma^\mu\psi$; $H_{NR}=-a_0$ at $m^0$, $-a_iv^i/m$ at $m^{-1}$; no spin-dependent leading term \\
$b_\mu$ & Vector & Odd & Odd & $\mathcal L_b=-b_\mu\,\bar\psi\,\gamma_5\gamma^\mu\,\psi$ \\
$b_i$ & Spatial part of $b_\mu$ & Odd & Odd & Generates $-\mathbf b\cdot\boldsymbol\sigma$ (matter) \\
$b_0$ & Temporal part of $b_\mu$ & Odd & Odd & Generates $+b_0(\boldsymbol\sigma\cdot\mathbf p)/m$ \\
$c_{\mu\nu}$ & Tensor, kinetic sector & Even & Odd & $\mathcal L_c=-c_{\mu\nu}\bar\psi\gamma^\mu D^\nu\psi$; direction-dependent kinetic modification; possible subleading $V_2$-type role via $g_sg_s$ \\
$H_{\mu\nu}$ & Antisymmetric tensor & Even & Odd & $\mathcal L_H=-\tfrac12 H_{\mu\nu}\,\bar\psi\,\sigma^{\mu\nu}\,\psi$ \\
$H_{ij}$ & Spatial part of $H_{\mu\nu}$ & Even & Odd & Magnetic-like; generates $-\boldsymbol{\mathcal H}_B\cdot\boldsymbol\sigma$ \\
$H_{0i}$ & Mixed part of $H_{\mu\nu}$ & Even & Odd & Electric-like; generates $-\tfrac1m\boldsymbol\sigma\cdot(\mathbf p\times\mathbf H_E)$ \\
$d_{\mu\nu}$ & Traceless tensor, kinetic sector & Even & Odd & $\mathcal L_d=\tfrac12 i\,\bar\psi\,d^{\mu\nu}\gamma_5\gamma_\mu\,\overleftrightarrow\partial_\nu\,\psi$ \\
$d_{i0}$ & Mixed part of $d_{\mu\nu}$ & Even & Odd & Generates $+d_{i0}\,m\,\sigma^i$ ($m^1$ enhancement) \\
$d_{ij}$ & Spatial part of $d_{\mu\nu}$ & Even & Odd & Generates $+d_{ij}\,p^j\,\sigma^i$ (velocity-dependent) \\
$d_{00}$ & Temporal part of $d_{\mu\nu}$ & Even & Odd & Generates $+d_{00}\,(\boldsymbol\sigma\cdot\mathbf p)$; nonzero because $d_{\mu\nu}$ is traceless, not antisymmetric \\
$e_\mu,\,f_\mu$ & Dimension-5 vectors & Odd & Odd & Velocity-dependent; contribute to $V_8$ and higher \\
$g_{\lambda\mu\nu}$ & Dimension-5 tensor & Odd & Odd & Listed for completeness in Ch.~2; not treated further \\
\bottomrule
\end{tabular}
\caption{SME fermion-sector coefficients used in this thesis.}
\label{tab:app-sme}
\end{table}

\paragraph{Operator-ordering convention.} The $b_\mu$ term's operator
ordering matters: this thesis uses $\gamma_5\gamma^\mu$ throughout (Ch.~1
Eq.~2, Ch.~2 Eq.~sme-minimal), not $\gamma^\mu\gamma_5$. Since
$\{\gamma_5,\gamma^\mu\}=0$, the two orderings are exact negatives of each
other. The companion symbolic notebook
\texttt{derivations/sympy/FW\_bmu\_term.ipynb} implements the opposite
ordering internally and consequently reports the opposite overall sign for
$H_{NR}(b_i)$ — a labelled convention choice, not an independent result;
see the notebook's own convention note for the explicit correspondence.

\section{Dobrescu--Mocioiu Coupling Constants and Potentials}

\begin{table}[h]
\centering
\begin{tabular}{@{}ll@{}}
\toprule
Symbol & Meaning \\
\midrule
$g_s$ & Scalar coupling constant \\
$g_p$ & Pseudoscalar coupling constant \\
$g_V$ & Vector coupling constant \\
$g_A$ & Axial-vector coupling constant \\
\bottomrule
\end{tabular}
\caption{Dobrescu--Mocioiu coupling constants (Ch.~1--2).}
\label{tab:app-dm-couplings}
\end{table}

Chapter~2 gives the complete closed-form catalogue
$V(\mathbf r)=\sum_{i=1}^{16}V_i(\mathbf r)$, transcribed from Cong et al.\
(2025). Table~\ref{tab:app-dm-potentials} lists the four potentials actually
derived from SME coefficients in Chapter~3, together with the
monopole--dipole example worked in Chapter~1.

\begin{table}[h]
\centering
\small
\begin{tabular}{@{}l>{\raggedright\arraybackslash}p{7.5cm}l@{}}
\toprule
Potential & Closed form & SME source \\
\midrule
$V_2$ & $-g_A^{(1)}g_A^{(2)}\,\tfrac{1}{4\pi}\,(\boldsymbol\sigma_1\cdot\boldsymbol\sigma_2)\,\tfrac{e^{-r/\lambda}}{r}$ & $b_i$, $d_{i0}$ \\
$V_3$ & $\tfrac{g^2}{4\pi}\big[\boldsymbol\sigma_1\cdot\boldsymbol\sigma_2-3(\boldsymbol\sigma_1\cdot\hat{\mathbf r})(\boldsymbol\sigma_2\cdot\hat{\mathbf r})\big]\tfrac{e^{-r/\lambda}}{r}$ & $H_{ij}$ \\
$V_{6,7}$ & $-\tfrac{1}{2mr^2}\big[(\boldsymbol\sigma_1\cdot\mathbf v)(\boldsymbol\sigma_2\cdot\hat{\mathbf r})\pm(\boldsymbol\sigma_1\cdot\hat{\mathbf r})(\boldsymbol\sigma_2\cdot\mathbf v)\big](1-r\,d/dr)\,y(r)$ & $H_{0i}$, $d_{ij}$ \\
$V_8$ & $\tfrac1r(\boldsymbol\sigma_1\cdot\mathbf v)(\boldsymbol\sigma_2\cdot\mathbf v)\,y(r)$ & $b_0$, $d_{ij}$, $d_{00}$ \\
$V_{9,10}$ & $-\tfrac{g_pg_s}{8\pi m}(\boldsymbol\sigma\cdot\hat{\mathbf r})\left(\tfrac1{r^2}+\tfrac1{\lambda r}\right)e^{-r/\lambda}$ (Ch.~1, one-spin example) & Not derived from an SME coefficient in Ch.~3; a database category in Ch.~4 \\
\bottomrule
\end{tabular}
\caption{Dobrescu--Mocioiu potentials derived or referenced in this thesis, with $y(r)=e^{-r/\lambda}/(4\pi)$.}
\label{tab:app-dm-potentials}
\end{table}

The per-coefficient matching sections in Chapter~3 (and the corresponding
theory notes in \texttt{docs/theory\_notes/}) also write $V_7$ and $V_8$ in
an informal, single-vertex form, for example an anticommutator
$\{(\boldsymbol\sigma\cdot\mathbf v),\,e^{-r/\lambda}/r\}$ for $V_8$, to
indicate which coefficient generates which potential class. The closed
two-body forms in Table~\ref{tab:app-dm-potentials} are the ones to cite
when the exact functional form is required.

Several further potential labels appear only in the Chapter~4 database, not
as a direct Chapter~3 derivation: $V_1$ (monopole--monopole), $V_{1a}$,
$V_{2+3}$, $V_{4+5}$, $V_{9+10}$, and $V_{12+13}$, combined labels used
because the underlying experiments do not distinguish the adjacent
catalogue potentials individually. $V_6$ and $V_{14}$ have no classified
dataset in the compiled database at all.

\section{Operators and Hamiltonians}

\begin{table}[h]
\centering
\begin{tabular}{@{}l>{\raggedright\arraybackslash}p{4cm}>{\raggedright\arraybackslash}p{6cm}@{}}
\toprule
Symbol & Meaning & Expression \\
\midrule
$H_0$ & Free Dirac Hamiltonian & $H_0=\boldsymbol\alpha\cdot\mathbf p+\beta m$ \\
$H$ & Full Hamiltonian & $H=\beta m+\varepsilon_{\text{even}}+\varepsilon_{\text{odd}}$ (Ch.~1's introductory sketch writes the equivalent $H=\boldsymbol\alpha\cdot\mathbf p+\beta m+\delta H_{\mathrm{SME}}$) \\
$H_{NR}$ & Non-relativistic Hamiltonian & After the Foldy--Wouthuysen transformation \\
$H_b$ & $b_\mu$ contribution & $b_\mu\gamma^0\gamma_5\gamma^\mu$ \\
$H_d$ & $d_{\mu\nu}$ contribution & $-d^{\mu\nu}p_\nu\,\Gamma_\mu$ \\
$\Gamma_\mu$ & $\gamma^0\gamma_5\gamma_\mu$ & $\Gamma_0=-\gamma_5$ (odd), $\Gamma_i=+\Sigma^i$ (even) \\
$\varepsilon_{\text{even}}$ & Even operator & Commutes with $\beta$ \\
$\varepsilon_{\text{odd}}$ & Odd operator & Anticommutes with $\beta$ \\
$U_{FW}$ & FW generator & $U_{FW}=\exp\!\left(\beta\varepsilon_{\text{odd}}/(2m)\right)$ \\
N/A & Master $1/m$ FW expansion & $H'\approx\beta m+\varepsilon_{\text{even}}+\dfrac{\beta\varepsilon_{\text{odd}}^2}{2m}-\dfrac{[\varepsilon_{\text{odd}},\varepsilon_{\text{even}}]}{4m}+O(m^{-2})$ \\
\bottomrule
\end{tabular}
\caption{Operators and Hamiltonians used in the Foldy--Wouthuysen reduction (Ch.~1--3).}
\label{tab:app-operators}
\end{table}

\section{Physical Quantities}

\begin{table}[h]
\centering
\begin{tabular}{@{}l>{\raggedright\arraybackslash}p{9.5cm}@{}}
\toprule
Symbol & Meaning \\
\midrule
$m$ & Fermion mass \\
$m_\phi$ & Mediator (boson) mass \\
$\lambda$ & Interaction range, $\lambda=\hbar/(m_\phi c)=1/m_\phi$ in natural units \\
$\mathbf p$ & Momentum operator, $-i\nabla$ \\
$\mathbf v$ & Velocity: single-particle NR relation $\mathbf v=\mathbf p/m$; in the two-body DM potentials, $\mathbf v=\mathbf v_{\text{rel}}$, the relative velocity of the pair \\
$\hat{\mathbf r}$ & Unit separation vector \\
$r$ & Distance between particles \\
$E$ & Energy eigenvalue (leading order: $m$) \\
$\langle\sigma_i\rangle_{\text{source}}$ & Macroscopic spin expectation of a polarised source, used in $b_i\leftrightarrow(g_A/2)\langle\sigma_i\rangle_{\text{source}}$ (Ch.~3) \\
\bottomrule
\end{tabular}
\caption{Physical quantities used throughout the thesis.}
\label{tab:app-physical}
\end{table}

\section{The Asymmetry Parameter}
\label{app:asymmetry}

The CPT asymmetry parameter is defined, for interaction channel
$\alpha\in\{s,p,V,A\}$, as
\begin{equation}
  A_\alpha = \frac{g_\alpha^f-g_\alpha^{\bar f}}{g_\alpha^f+g_\alpha^{\bar f}},
\end{equation}
where $g_\alpha^f$ is the effective coupling constant for fermion $f$ in
channel $\alpha$ and $g_\alpha^{\bar f}$ is the corresponding coupling for
its antifermion; the channel index is dropped when clear from context (as
in Chapter~3's coefficient-by-coefficient derivations), and Chapter~4
promotes the same quantity to a function of the interaction range,
$A_\alpha(\lambda)$, evaluated on a grid for each matched dataset pair.

Three properties of $A_\alpha$ hold throughout the thesis, independent of
which coefficient is being discussed:
\begin{itemize}
  \item Under exact CPT symmetry, $A_\alpha=0$ for every channel.
  \item For a genuinely CPT-odd, exactly signed coupling
    ($g_\alpha^{\bar f}=-g_\alpha^f$), the denominator of $A_\alpha$
    vanishes identically, so $A_\alpha$ \emph{diverges} rather than
    saturating at $1$ (Ch.~3, confirmed with a symbolic limit calculation).
  \item SPINDEP's actual inputs, $g_\alpha^f$ and $g_\alpha^{\bar f}$, are
    independent, always-positive experimental \emph{upper bounds}, not
    signed measured couplings. A ratio of two independent positive bounds
    of very different tightness approaches $\pm1$ whenever one bound is far
    tighter than the other, regardless of the true CPT parity of the
    underlying physics, a sensitivity-gap effect, not a signature of any
    particular coefficient (Ch.~1, 3, 4).
\end{itemize}
A further, distinct point established in Chapter~4: statistical
significance is not the same claim as evidence for CPT violation. Every one
of the ten matched matter--antimatter pairs compiled in Chapter~4 is
significant at effectively $p\to0$ after the autocorrelation correction
described there, yet none can currently be read as CPT-violation evidence
rather than a sensitivity gap between the matter- and antimatter-sector
measurements being compared.

\section{Statistical and Database Notation}

\begin{table}[h]
\centering
\begin{tabular}{@{}l>{\raggedright\arraybackslash}p{9.5cm}@{}}
\toprule
Symbol & Meaning \\
\midrule
\texttt{\{V\}\{Author\}\{Year\}\_m\_abs\_\{sector\}.csv} & Filename convention for a compiled constraint dataset: potential label, source, and fermion sector encoded in the name \\
\texttt{ee}, \texttt{eebar}, \texttt{ep}, \texttt{epbar}, \texttt{emu}, \texttt{emubar}, \ldots & Fermion-sector labels; \texttt{bar}-suffixed sectors are antimatter (e.g.\ \texttt{eebar} = electron--positron) \\
\texttt{UNKNOWN} & Potential label the parser could not resolve from a filename; retained rather than dropped or guessed, so the classification gap stays visible \\
$g_Ag_A$, $g_pg_s$, $g_sg_s$, $g_Vg_V$, $g_Ag_V$, $g_pg_p$ & Coupling-family classification tags for a dataset (which coupling type each side of the interaction uses), not a literal numeric product \\
$\chi^2$ & Chi-squared statistic, computed both with a uniform 10\% fractional uncertainty and with a curvature-weighted per-point uncertainty \\
$\mathrm{dof}_{\text{eff}}$ & Effective degrees of freedom, corrected for autocorrelation among the 300 interpolated grid points (found to be 6--21, versus the naive 300) \\
$p_{\text{eff}}$ & Significance recomputed against $\mathrm{dof}_{\text{eff}}$ \\
$\overline{|A_\alpha|}$ & Mean absolute asymmetry across the grid for one matched pair \\
95\% CI & Bootstrap (2000 resamples) confidence interval on $\overline{|A_\alpha|}$ \\
\bottomrule
\end{tabular}
\caption{Statistical and database notation used in Chapter~4.}
\label{tab:app-stats}
\end{table}

\section{Key Signs}

\begin{table}[h]
\centering
\small
\begin{tabular}{@{}l>{\raggedright\arraybackslash}p{6.5cm}@{}}
\toprule
Expression & Meaning \\
\midrule
$-\mathbf b\cdot\boldsymbol\sigma$ & Matter NR Hamiltonian for $b_i$, from $\gamma^0\gamma_5\gamma^i=-\Sigma^i$ \\
$+\mathbf b\cdot\boldsymbol\sigma$ & Antimatter NR Hamiltonian for $b_i$, from charge conjugation (CPT-odd) \\
$+b_0(\boldsymbol\sigma\cdot\mathbf p)/m$ & NR Hamiltonian for $b_0$; an odd operator becomes even via FW at $m^{-1}$ \\
$-\boldsymbol{\mathcal H}_B\cdot\boldsymbol\sigma$ & NR Hamiltonian for $H_{ij}$; standard Zeeman coupling, CPT-even \\
$-\tfrac1m\boldsymbol\sigma\cdot(\mathbf p\times\mathbf H_E)$ & NR Hamiltonian for $H_{0i}$; odd operator at $m^{-1}$ via FW \\
$+d_{i0}\,m\,\sigma^i$ & NR Hamiltonian for $d_{i0}$; matches Kostelecký \& Lane's Eq.~(4) exactly, sign included \\
$+d_{ij}\,p^j\,\sigma^i$ & NR Hamiltonian for $d_{ij}$; structure confirmed, overall sign open \\
$+d_{00}\,(\boldsymbol\sigma\cdot\mathbf p)$ & NR Hamiltonian for $d_{00}$; structure confirmed, overall sign open \\
$A_\alpha\to\infty$ & Formal result for an exact signed CPT-odd relation; denominator vanishes \\
$A_\alpha\to\pm1$ & Actual result from one-sided experimental bounds; a sensitivity-gap effect, not CPT evidence \\
\bottomrule
\end{tabular}
\caption{Key sign results and their origin.}
\label{tab:app-signs}
\end{table}

\section{Commutation and Anticommutation Relations}

\begin{table}[h]
\centering
\begin{tabular}{@{}l>{\raggedright\arraybackslash}p{7cm}@{}}
\toprule
Relation & Meaning \\
\midrule
$\{\gamma_5,\gamma^\mu\}=0$ & $\gamma_5$ anticommutes with every gamma matrix \\
$\gamma^0\gamma_5\gamma^i=-\Sigma^i$ & The spatial $b_i$ term is even (block-diagonal) \\
$\gamma^0\gamma_5\gamma^0=\gamma_5$ & The temporal $b_0$ term is odd (off-diagonal) \\
$\gamma^0\sigma^{ij}=\varepsilon^{ijk}\Sigma^k$ & The $H_{ij}$ term is even \\
$\gamma^0\sigma^{0i}=-\alpha^i$ & The $H_{0i}$ term is odd \\
$\Gamma_0=-\gamma_5$ & The $d_{\mu\nu}$ temporal piece is odd \\
$\Gamma_i=+\Sigma^i$ & The $d_{\mu\nu}$ spatial piece is even \\
$[\varepsilon_{\text{even}},\beta]=0$ & Definition of an even operator \\
$\{\varepsilon_{\text{odd}},\beta\}=0$ & Definition of an odd operator \\
\bottomrule
\end{tabular}
\caption{Commutation and anticommutation relations underlying the even/odd decomposition.}
\label{tab:app-commutation}
\end{table}

\section{CPT Classification}
\label{app:cpt-table}

Table~\ref{tab:app-cpt} summarises how each coefficient's CPT parity is
established. Most follow directly from the bilinear transformation rule
$\mathrm{CPT}[\bar\psi\Gamma\psi]=\eta_\Gamma\bar\psi\Gamma\psi$ given in
Section~\ref{tab:app-dirac}. Two cases do not, and are worth remembering
explicitly: $d_{\mu\nu}$'s operator is $\gamma_5\gamma_\mu$ paired with a
derivative $\partial_\nu$, and the derivative itself carries a CPT
transformation the plain bilinear rule does not capture; and $a_\mu$'s
bilinear alone ($\gamma^\mu$, $\eta_\Gamma=+1$) would naively suggest
CPT-even, yet $a_\mu$ is classified CPT-odd in the source literature. Both
are therefore stated directly, following Kostelecký \& Lane, rather than
derived from the bilinear rule.

\begin{table}[h]
\centering
\small
\begin{tabular}{@{}l>{\raggedright\arraybackslash}p{7cm}l@{}}
\toprule
Coefficient & Basis for CPT parity & Result \\
\midrule
$b_\mu$ & Bilinear rule: $\mathrm{CPT}[\bar\psi\gamma_5\gamma^\mu\psi]=-\bar\psi\gamma_5\gamma^\mu\psi$ & CPT-odd \\
$H_{\mu\nu}$ & Bilinear rule: $\mathrm{CPT}[\bar\psi\sigma^{\mu\nu}\psi]=+\bar\psi\sigma^{\mu\nu}\psi$ & CPT-even \\
$c_{\mu\nu}$ & Bilinear rule; the paired derivative does not change the vector current's parity & CPT-even \\
$d_{\mu\nu}$ & Stated directly (Kostelecký \& Lane); not derivable from the bilinear rule alone, since the operator is derivative-coupled & CPT-even \\
$a_\mu,\,e_\mu,\,f_\mu,\,g_{\lambda\mu\nu}$ & Stated directly (Kostelecký \& Lane) & CPT-odd \\
$b_\mu$ (charge conj.) & Current is $C$-even, but the coefficient itself flips sign under $C$ & Matter couples as $-b_\mu$, antimatter as $+b_\mu$ \\
$H_{\mu\nu}$, $d_{\mu\nu}$ (charge conj.) & Current is $C$-even and the coefficient does not flip & Same sign for matter and antimatter \\
\bottomrule
\end{tabular}
\caption{CPT classification of the SME coefficients treated in this thesis.}
\label{tab:app-cpt}
\end{table}

\section{Summary of Key Results}

\begin{enumerate}
  \item \textbf{$b_\mu$ sign flip.} Matter couples as $-\mathbf b\cdot\boldsymbol\sigma$; antimatter as $+\mathbf b\cdot\boldsymbol\sigma$: a genuine CPT-odd signature.
  \item \textbf{$H_{\mu\nu}$, $d_{\mu\nu}$, $c_{\mu\nu}$ same sign.} Matter and antimatter couple identically (CPT-even). $d_{\mu\nu}$'s evenness follows Kostelecký \& Lane's direct classification rather than the bilinear rule.
  \item \textbf{$A_\alpha$ ambiguity.} A ratio of one-sided upper bounds gives $|A_\alpha|\to1$ from a sensitivity gap alone, independent of the true CPT parity, confirmed both formally (Ch.~1, 3) and empirically across every matched pair in the compiled database (Ch.~4), including the CPT-even $g_sg_s$ pair.
  \item \textbf{Statistical significance is not CPT evidence.} Every matched pair in Chapter~4 is significant after the autocorrelation correction, but significance establishes that an asymmetry is real, not that it is caused by CPT violation.
  \item \textbf{$d_{i0}$ mass enhancement.} $H_{NR}(d_{i0})=+d_{i0}\,m\,\sigma^i$ matches Kostelecký \& Lane's Eq.~(4) exactly, sign included.
  \item \textbf{$d_{ij}$, $d_{00}$ sign.} Operator structure is correct; the overall sign remains open pending a wavefunction-renormalization treatment of the kinetic sector.
  \item \textbf{The central empirical limitation.} Of 273 compiled datasets, only 19 touch an antimatter sector, concentrated in the $e$-$\bar p$, $e$-$e^+$, and $e$-$\bar\mu$ channels, the reason a meaningful matter--antimatter comparison is currently possible in only ten matched pairs.
\end{enumerate}

\end{document}
